\documentclass[a4paper,14pt]{extarticle}

\usepackage{styles}

\title{Technical Report: GPU-Accelerated Genetic Algorithms\\Using CUDA Optimization for XOR Perceptron and TSP Benchmarks}
\author{Aleksieienko A.\\
{\small National University of Kyiv-Mohyla Academy, Department of Computer Science}}
\date{}

\begin{document}
\ifdefined\reportcomposed\else
\maketitle
\thispagestyle{empty}
\fi

\subsection*{Disclaimer}

\textbf{This is an AI-generated technical report and is not part of the thesis.}

This document was generated by Claude Opus 4.5 (Anthropic) on January~30,
2026, and subsequently edited for accuracy. It provides a technical overview of the codebase
implementation as of the report generation date.

\textbf{Purpose:} This report serves as supplementary documentation describing the software
architecture, implementation details, and technical specifications of the research project.

\textbf{Source Code:} The complete codebase is available at:\\
\url{https://github.com/ukma-morphonas-lab/py-cuda-ga}

For the actual thesis text please refer to the main thesis document.

\textbf{Date:} January 30, 2026\\
\textbf{Author:} Aleksieienko A.


%%% ===================================================================
\section{Project Overview}
%%% ===================================================================

\subsection{Purpose}

This project implements a GPU-accelerated genetic algorithm (GA) research framework using
NVIDIA CUDA via Numba JIT compilation and CuPy. It benchmarks progressive CUDA optimization
stages for two GA benchmark problems---XOR perceptron weight training and the Traveling
Salesman Problem (TSP)---on entry-level mobile GPU hardware (NVIDIA RTX~2050, Compute
Capability~8.6).

Since each individual is evaluated independently, fitness computation constitutes an
embarrassingly parallel workload well suited to GPU acceleration~\cite{cheng2020parallel}.
Prior CUDA-based GA implementations report speedups ranging from $1.5\times$ to over
$20\times$ on dedicated compute GPUs~\cite{plichta2014implementation}, but this project
focuses on examining entry-level hardware and quantifying transfer overhead.

\subsection{Research Question}

What GPU speedup is achievable for genetic algorithm fitness evaluation on entry-level mobile
hardware, and what is the impact of host--device transfer overhead across incremental CUDA
optimization stages?

\subsection{Methodology}

\begin{itemize}
  \item \textbf{GPU Framework:} Numba CUDA JIT compilation with CuPy for device management
  \item \textbf{Hardware:} NVIDIA RTX~2050 (16~SMs, 3~GB VRAM, Compute Capability~8.6)
  \item \textbf{Benchmarks:} XOR perceptron (continuous optimization) and 50-city TSP
        (combinatorial optimization)
  \item \textbf{Optimization Approach:} Five incremental GPU stages applied progressively
        (fitness kernel $\to$ shared memory $\to$ device-resident population $\to$ pinned
        host memory $\to$ GPU-side RNG)
  \item \textbf{Experimental Design:} Paired CPU/GPU runs with identical random seeds
  \item \textbf{Statistical Validation:} Paired $t$-tests, Wilcoxon signed-rank tests,
        bootstrap confidence intervals, power analysis
\end{itemize}


%%% ===================================================================
\section{Source Code Architecture}
%%% ===================================================================

\subsection{Core Modules and Responsibilities}

\paragraph{CUDA Configuration Module (\texttt{cuda/config/})}

Provides GPU device query, thread/block configuration calculation, and occupancy estimation.

\begin{itemize}
  \item \textbf{optimal\_config\_getter.py}: GPU properties query via CuPy, optimal
        threads-per-block and blocks-per-grid calculation, occupancy estimation for
        Compute Capability~8.6
\end{itemize}

Key classes and functions:
\begin{itemize}
  \item \texttt{GPUProperties} --- dataclass encapsulating device properties (compute
        capability, memory, clock rates, thread limits)
  \item \texttt{calculate\_occupancy()} --- computes GPU occupancy percentage
  \item \texttt{get\_threads\_per\_block()} / \texttt{get\_blocks\_per\_grid()} ---
        calculate optimal launch configuration based on population size
  \item \texttt{calculate\_optimal\_threads\_config()} --- main configuration calculator
\end{itemize}

\paragraph{CUDA Profiling Module (\texttt{cuda/profiling/})}

Captures GPU state snapshots for test instrumentation.

\begin{itemize}
  \item \textbf{gpu\_state\_snapshot.py}: \texttt{GPUStateSnapshot} class that captures
        device count, memory utilization, and device name via CuPy memory pool integration;
        provides JSON-serializable snapshots used by the test logger
\end{itemize}

\paragraph{GA Test Infrastructure (\texttt{genetic/tests/})}

Shared testing framework for all GA benchmarks.

\begin{itemize}
  \item \textbf{constants.py}: Global \texttt{RANDOM\_SEED = 42} for reproducibility
  \item \textbf{logger.py}: \texttt{CUDA\_GA\_TestLogger} class for comprehensive test
        instrumentation---environment info, timing, parameters, and results saved to JSON
  \item \textbf{fitness\_metrics.py}: Abstract \texttt{FitnessMetric} base class with
        \texttt{MaximizationMetric} and \texttt{MinimizationMetric} implementations
  \item \textbf{statistical\_analysis\_base.py}: \texttt{StatisticalResults} dataclass and
        \texttt{compare\_cpu\_gpu\_results()} function implementing paired $t$-tests,
        Wilcoxon signed-rank tests, and 10\,000-sample bootstrap confidence intervals
  \item \textbf{visualization\_base.py}: Base class for result visualizations
\end{itemize}

\paragraph{Perceptron Tests (\texttt{genetic/tests/perceptron/})}

Progressive GPU optimization benchmarks for XOR perceptron weight training.

\begin{itemize}
  \item \textbf{test\_xor\_perceptron\_v0\_cpu.py} (v0 --- CPU baseline): NumPy-based serial
        evaluation with \texttt{PerceptronXOR} and \texttt{GA\_PerceptronXOR} classes
  \item \textbf{test\_xor\_perceptron\_v1.py} (v1 --- GPU fitness kernel):
        \texttt{@cuda.jit} kernel, 256~threads/block, one thread per individual
  \item \textbf{test\_xor\_perceptron\_v2.py} (v2 --- + shared memory): XOR input/output
        patterns cached in shared memory
  \item \textbf{test\_xor\_perceptron\_v3.py} (v3 --- + device-resident population):
        persistent GPU buffers with device-side selection and merge kernels
  \item \textbf{test\_xor\_perceptron\_v4.py} (v4 --- + pinned host memory): page-locked
        arrays for remaining transfers
  \item \textbf{test\_xor\_perceptron\_v5.py} (v5 --- + GPU RNG): Xoroshiro128p parallel
        PRNG on device
\end{itemize}

\paragraph{TSP Tests (\texttt{genetic/tests/salesman/})}

Progressive GPU optimization benchmarks for 50-city Traveling Salesman Problem.

\begin{itemize}
  \item \textbf{test\_salesman\_simple\_ga.py} (v1 --- CPU baseline): \texttt{Graph} class
        with Euclidean distances, NumPy-based population management
  \item \textbf{test\_salesman\_simple\_ga\_v2.py} (v2 --- GPU fitness kernel):
        \texttt{calculate\_fitness\_kernel()} for parallel path cost evaluation
  \item \textbf{test\_salesman\_simple\_ga\_v3.py} (v3 --- + device-resident population):
        further GPU optimizations
  \item \textbf{test\_salesman\_simple\_ga\_v4.py} (v4 --- + GPU RNG, pinned memory,
        pre-allocated buffers): most optimized version
\end{itemize}

\paragraph{CUDA Unit Tests (\texttt{cuda/tests/})}

Initial CUDA kernel tests, ran to explore the setup before the genetic algorithms implementation.

\begin{itemize}
  \item \textbf{test\_vector\_addition.py}: Basic CUDA vector addition kernel
  \item \textbf{test\_matmul.py} / \textbf{test\_matmul\_v2.py}: Matrix multiplication
        kernels with progressive optimization
  \item \textbf{test\_cupy\_matmul\_v3.py}: CuPy-based matrix multiplication
\end{itemize}


\subsection{Data Flow}

\begin{verbatim}
Test Invocation (pytest)
    |
    v
Session-scoped Logger (CUDA_GA_TestLogger)
    |  - captures environment info (GPU, CPU, memory)
    |  - starts timing
    v
GA Execution (CPU or GPU version)
    |
    v
For each generation:
    1. Evaluate fitness (parallel on GPU or serial on CPU)
       - GPU: @cuda.jit kernel, one thread per individual
       - CPU: NumPy vectorized evaluation
    2. Selection (tournament or roulette)
    3. Crossover + Mutation
       - GPU v5+: device-side RNG (Xoroshiro128p)
       - Earlier versions: host-generated random data
    4. Transfer results (GPU versions only)
       - Pinned memory (v4+) or pageable transfers
    |
    v
Results Collection
    |  - wall-clock time (CPU & GPU)
    |  - best fitness / path cost
    |  - transfer overhead percentage
    v
Statistical Analysis (compare_cpu_gpu_results)
    |  - paired t-test (time)
    |  - Wilcoxon signed-rank (fitness)
    |  - bootstrap speedup CI (10,000 samples)
    v
JSON Output (logs/results/)
\end{verbatim}


%%% ===================================================================
\section{Key Parameters and Metrics}
%%% ===================================================================

\subsection{XOR Perceptron Benchmark Parameters}

\begin{itemize}
  \item \textbf{Chromosome encoding:} two real-valued weights
  \item \textbf{Fitness:} correct classifications out of four XOR input patterns (theoretical
        maximum~2 for a single-layer perceptron, since XOR is linearly inseparable)
  \item \textbf{Population:} 100 individuals
  \item \textbf{Generations:} 500 (fixed, no early termination)
  \item \textbf{Paired runs:} 30
\end{itemize}

The problem was chosen deliberately as a lightweight, deterministic benchmark for measuring
execution-time differences between CPU and GPU implementations; its linear inseparability is
acknowledged and irrelevant to the performance comparison, which concerns wall-clock time
rather than classification accuracy.

\subsection{TSP Benchmark Parameters}

\begin{itemize}
  \item \textbf{Chromosome encoding:} city permutation
  \item \textbf{Fitness:} total path cost (minimization)~\cite{hussain2017genetic}
  \item \textbf{Cities:} 50 (random graph, fixed seed)
  \item \textbf{Population:} 200 individuals
  \item \textbf{Generations:} 100 (fixed, no early termination)
  \item \textbf{Paired runs:} 50
\end{itemize}

Both problems execute a fixed number of generations regardless of solution quality, eliminating
early-termination variance in timing measurements.

\subsection{GPU Optimization Stages}

The perceptron benchmark progressed through six versions:\\
\textbf{v0 (CPU baseline):} NumPy-based serial evaluation, crossover, and mutation.\\
\textbf{v1 (GPU fitness kernel):} \texttt{@cuda.jit} kernel with 256 threads per block;
one thread evaluates one individual.\\
\textbf{v2 (+ shared memory):} XOR input/output patterns cached in shared memory,
reducing repeated global memory access~\cite{nvidia2011cuda}.\\
\textbf{v3 (+ device-resident population):} Persistent GPU buffers across generations
with device-side selection and merge kernels, reducing host--device transfers from
$\mathcal{O}(\text{generations})$ to $\mathcal{O}(1)$~\cite{harris2012optimize}.\\
\textbf{v4 (+ pinned host memory):} Page-locked arrays for remaining transfers,
enabling DMA without staging copies~\cite{nvidia2024cuda}.\\
\textbf{v5 (+ GPU RNG):} Xoroshiro128p parallel PRNG on device, eliminating
per-generation transfer of random data~\cite{xoroshiro}.

The TSP benchmark followed an analogous progression (v0--v4): CPU baseline, GPU fitness
kernel, device-resident population, and pinned memory with GPU RNG.

\subsection{Statistical Validation Parameters}

\begin{itemize}
  \item \textbf{Design:} Paired experimental --- GPU and CPU executions share identical seeds
        per run
  \item \textbf{Time comparison:} Paired $t$-test
  \item \textbf{Fitness comparison:} Wilcoxon signed-rank test
  \item \textbf{Speedup CI:} 95\% via 10\,000 bootstrap resamples
  \item \textbf{Significance threshold:} $\alpha = 0.05$
  \item \textbf{Power analysis:} $N = 30$ provides $> 89\%$ power; $N = 50$ provides
        $> 97\%$ power for large effects (Cohen's $d > 0.8$)~\cite{cohen1988power}
\end{itemize}

\subsection{GPU Hardware Configuration}

\begin{itemize}
  \item \textbf{Device:} NVIDIA RTX~2050 (mobile)
  \item \textbf{Compute Capability:} 8.6 (Ampere)
  \item \textbf{Streaming Multiprocessors:} 16
  \item \textbf{VRAM:} 3~GB
  \item \textbf{CUDA version:} 12.4
  \item \textbf{Numba version:} 0.61.2
\end{itemize}


%%% ===================================================================
\section{Experiment Result Files}
%%% ===================================================================

Test results are saved to \texttt{genetic/tests/\{perceptron,salesman\}/logs/} as structured
JSON files.

\textbf{Per test session:}
\begin{itemize}
  \item \textbf{results/*session*.json}: Complete session data including:
    \begin{itemize}
      \item Environment info (GPU device, CPU, system memory, CUDA version)
      \item Test parameters (population, generations, seed)
      \item Wall-clock timing for each run
      \item Best fitness / path cost per run
    \end{itemize}
\end{itemize}

\textbf{Statistical analysis output:}
\begin{itemize}
  \item Mean and standard deviation for GPU and CPU fitness
  \item Mean and standard deviation for GPU and CPU wall-clock time
  \item Paired $t$-test $p$-value (time) and Wilcoxon $p$-value (fitness)
  \item Speedup factor with 95\% bootstrap confidence interval
  \item Significance flags for both time and fitness comparisons
  \item Formatted text summary via \texttt{StatisticalResults.summary()}
\end{itemize}

\textbf{Profiling data} (\texttt{profiling/nsight/}):
\begin{itemize}
  \item NVIDIA Nsight Systems profiling output and analysis for TSP GPU kernels
\end{itemize}


%%% ===================================================================
\section{Technical Implementation Details}
%%% ===================================================================

\subsection{CUDA Kernel Design}

GPU fitness evaluation uses \texttt{@cuda.jit} decorated Numba kernels where one CUDA thread
evaluates one individual. Thread and block configuration is computed dynamically using
\texttt{optimal\_config\_getter.py} based on population size and device properties.

For the XOR perceptron, the kernel computes forward-pass outputs for all four input patterns
per thread. For TSP, the kernel sums Euclidean distances across the full city permutation per
thread.

\subsection{Memory Optimization}

Four memory optimization stages were applied:

\begin{enumerate}
  \item \textbf{Shared memory} (v2): Frequently accessed read-only data (XOR patterns)
        cached in per-block shared memory, reducing global memory bandwidth
        pressure~\cite{nvidia2011cuda}
  \item \textbf{Device-resident population} (v3): Population buffers persist on GPU across
        generations; selection and merge performed via device-side kernels, reducing
        host--device transfers from $\mathcal{O}(\text{generations})$ to
        $\mathcal{O}(1)$~\cite{harris2012optimize}
  \item \textbf{Pinned host memory} (v4): Page-locked arrays allocated via
        \texttt{cuda.pinned\_array()} for remaining transfers, enabling DMA without
        staging copies~\cite{nvidia2024cuda}
  \item \textbf{GPU-side RNG} (v5): Xoroshiro128p parallel PRNG initialized per-thread on
        device, eliminating per-generation transfer of random data~\cite{xoroshiro}
\end{enumerate}

\subsection{Reproducibility Guarantees}

\begin{itemize}
  \item Global \texttt{RANDOM\_SEED = 42} defined in \texttt{constants.py}
  \item Paired experimental design: identical seeds for CPU and GPU runs within each pair
  \item Session-scoped pytest fixtures capture full environment state (GPU properties, CPU
        info, system memory) at test start
  \item All results serialized to JSON for post-hoc analysis
\end{itemize}

\subsection{Code Quality}

\begin{itemize}
  \item Abstract base classes for fitness metrics (\texttt{FitnessMetric})
  \item Dataclass-based GPU properties (\texttt{GPUProperties})
  \item Session-scoped test instrumentation with automatic JSON export
  \item Separation of concerns: CUDA configuration, profiling, GA logic, statistical analysis,
        and visualization each in dedicated modules
  \item Progressive versioning (v0--v5) isolates the effect of each optimization stage
\end{itemize}


%%% ===================================================================
\section{Key Insights and Design Decisions}
%%% ===================================================================

\subsection{Progressive Optimization Strategy}

Each GPU optimization was applied incrementally rather than all at once. This versioned
approach (v0--v5) isolates the performance contribution of each technique and produces a clear
optimization narrative for the thesis. The final GPU version includes all prior optimizations.

\subsection{Paired Experimental Design}

GPU and CPU executions share identical random seeds per run, ensuring that any measured
difference in wall-clock time or solution quality arises from the execution platform rather
than stochastic variation in the evolutionary process.

\subsection{Entry-Level Hardware Focus}

The RTX~2050 (16~SMs, 3~GB VRAM) represents entry-level mobile GPU hardware. Results
demonstrate that meaningful GPU acceleration is achievable on such hardware with moderate
population sizes, though larger populations are needed to approach speedup factors reported
for high-end GPUs in prior work~\cite{plichta2014implementation}.

\subsection{Transfer Overhead as Primary Bottleneck}

Host--device transfer overhead consumed 25--29\% of GPU wall time across both benchmarks,
representing the primary performance bottleneck. Device-resident evolution (v3) and GPU-side
RNG (v5) were the most effective individual optimizations because they directly reduce transfer
volume.

\subsection{Statistical Rigor}

Multiple paired runs (30 for XOR, 50 for TSP) with power analysis confirmation ensure
statistical validity. Two complementary tests---parametric (paired $t$-test) and
non-parametric (Wilcoxon signed-rank)---guard against distributional assumptions.
Bootstrap confidence intervals for speedup provide uncertainty quantification beyond
point estimates~\cite{cohen1988power}.


%%% ===================================================================
\section{Dependencies}
%%% ===================================================================

\textbf{Core:}
\begin{itemize}
  \item \texttt{numba==0.61.2} (CUDA JIT compilation)
  \item \texttt{numpy>=1.20.0} (numerical computing)
  \item \texttt{cuda-python==13.0.1} (CUDA Python bindings)
  \item \texttt{cupy-cuda12x} (GPU array library and device management)
\end{itemize}

\textbf{Statistical Analysis:}
\begin{itemize}
  \item \texttt{scipy>=1.10.0} (paired $t$-tests, Wilcoxon, bootstrap)
\end{itemize}

\textbf{Visualization:}
\begin{itemize}
  \item \texttt{matplotlib==3.10.0} (plotting)
  \item \texttt{pandas==2.3.3} (data manipulation)
\end{itemize}

\textbf{Testing:}
\begin{itemize}
  \item \texttt{pytest>=8.1} (test framework)
  \item \texttt{pytest-benchmark==5.2.0} (performance benchmarking)
  \item \texttt{psutil==5.9.0} (system resource monitoring)
\end{itemize}

\textbf{Profiling:}
\begin{itemize}
  \item \texttt{snakeviz==2.2.2} (profiling visualization)
\end{itemize}


\ifdefined\reportcomposed\else
\newpage
\printbibliography
\fi

\end{document}
